% introduccion
\thispagestyle{empty}
\PageDecor{3}
\SectionTag{SECCIÓN INFORMATIVA}
\DocHeading{Introducción}
\DocPara{La intención de este cuadernillo es aportar información sobre distintos puntos de la vida estudiantil e institucional. Recorriendo distintos niveles que van desde las materias que vas a tener hasta cuestiones relacionadas a la cursada, como el descuento estudiantil en pasajes o la libreta de calificaciones.}
\DocPara{Como Centro de Estudiantes creemos que la autonomía estudiantil es un valor fundamental para la comunidad de este profesorado. Autonomía estudiantil significa que los y las estudiantes debatamos entre nosotros/as cuáles son nuestras necesidades y demandas, nos organicemos en pos de conquistarlas, nos fortalezcamos y así aportemos al crecimiento del Profesorado. Dicha autonomía rige la democracia estudiantil donde somos los y las estudiantes -y sólo nosotros- quienes decidimos cómo nos organizamos.}
\DocPara{La vida estudiantil tiene distintos aspectos. No se limita sólo a las clases, los exámenes y lo que demandan (cursar, sacar fotocopias, estudiar, etcétera), por más que sean los aspectos centrales. Alrededor de ese fin que es recibir una formación adecuada y recibirse como docente, hay muchas realidades. Llegar al profesorado, juntarse a estudiar con compañeros/as, asistir a espacios de formación alternativos, o participar de discusiones sobre las condiciones de estudio y nuestros derechos. Al mismo tiempo la realidad de cada estudiante es diferente a la de otro. Como Centro pretendemos trabajar con todas esas realidades, creando junto a cada estudiante herramientas y espacios que colaboren en hacer más fácil la permanencia en la carrera, articulando además con los otros claustros de la institución (el docente, el no docente y el graduado).}
\DocPara{Si estás ingresando por primera vez a estudiar alguna de las tres carreras que ofrece la institución, te damos la bienvenida y nos ponemos a tu disposición para orientarte en esta etapa.}
\DocPara{Si ya estudias una carrera, te ofrecemos esta información que puede llegar a serte útil y esperamos podamos trabajar juntos en este ciclo lectivo.}
